\section{Introduction}
\label{sec:intro}

Multi-agent debate (MAD) systems built on large language models (LLMs) have shown strong performance on complex reasoning~\citep{du2024improving,li2024improving,meng2026small,wang2025learning}, embodied action~\citep{shridhar2020alfworld,wang2025madra}, and planning~\citep{hu2025debate,li2025swe,ma2024agentboard} tasks, where agents iteratively discuss, critique, and refine each other's outputs~\citep{chan2023chateval,liang2024encouraging,more2026theramind}. To support interactions beyond a single round, recent MAD systems add a \emph{shared memory} that persists intermediate reasoning, past actions, and episodic trajectories across rounds~\citep{anokhin2024arigraph,wang2023voyager,zhang2025g,zhong2024memorybank}.

While shared memory boosts long-horizon reasoning, it also opens a critical vulnerability: a corrupted memory state, which can subsequently contaminate all downstream memory-augmented reasoning \citep{torra2026memory}. The corruption arises because individual agents are inherently imperfect: LLMs hallucinate, agree sycophantically ~\citep{banks2026ai,carro2024flattering,sharma2023towards}, or confidently assert incorrect claims~\citep{shinn2023reflexion,yang2025eligibility,zhang2025g}. These failures do not cancel under debate. Across recent works, three failure patterns present (Figure \ref{fig:intro}): (i) an over-confident contributor pushes a hallucinated entry past hedged auditors who defer to its confidence rather than challenge it ~\citep{du2024improving,wang2024rethinking,xiong2023examining}, producing a corrupted memory that reads like established fact; (ii) over-confident auditors veto a tentative but correct contribution \citep{cui2024or,zhang2025falsereject}, producing an over-curated memory that drops truly useful facts; (iii) all agents hedge (or over-confident) and nothing certain is committed \citep{song2025llms,tang2026value,tomani2024uncertainty}, leaving memory under-populated. In all three cases, the commit decision depends on agents' self-reported confidence rather than on any check against the memory state itself, so debate alone cannot filter the errors.

% While shared memory boosts long-horizon reasoning, it also opens a critical vulnerability: a corrupted entry written by one agent can contaminate all downstream memory-augmented reasoning~\citep{torra2026memory}. The corruption has two sources. First, early reasoning steps are produced without full task context and often contain inaccurate or outdated conclusions; once written to memory, later agents retrieve them as established facts and the errors propagate~\citep{shinn2023reflexion,zhang2025g}. Second, individual agents are inherently imperfect: LLMs hallucinate, agree sycophantically~\citep{sharma2023towards}, and confidently assert incorrect claims~\citep{fanous2025syceval}, which degrades the trustworthiness of contributed memory entries even without external adversaries.  One might expect debate itself to dilute such errors, but recent work shows the opposite: a single confidently wrong agent can pull the group toward an incorrect consensus~\citep{du2024improving,wang2024rethinking}, and agents tend to defer to early assertions rather than evaluate them independently~\citep{xiong2023examining}.


Existing safeguards either filter proposed entries via heuristics or auxiliary LLM classifiers~\citep{alon2023detecting,wang2024self}, or rely on consensus-based validation~\citep{wei2025memguard}. Both deliver real protection, but still produce calibration through AI models that carry the same risks of hallucination and sycophancy. More importantly, these methods are designed for a single agent reading its own memory, and do not model multiple agents continuously writing to (and potentially colluding over) a shared memory space. This gap motivates a critical question: \emph{how can we calibrate memory updates in a MAD system without relying on any single agent's judgment?}

\begin{wrapfigure}{r}{0.62\textwidth}
  \vspace{-22pt}
  \centering
  \includegraphics[width=0.62\textwidth]{figures/intro.pdf}
  \vspace{-13pt}
  \caption{Failures of debate-driven memory commitment.}
  \vspace{-5pt}
  \label{fig:intro}
\end{wrapfigure}
Answering this question requires rethinking the cross-agent dynamics of debate itself. Since MAD agents already inspect and critique each other's contributions, we formulate their interaction with memory as a \textbf{zero-trust memory game}. In this formulation, the agent proposing an update acts as a \emph{contributor} in each debating round, while the remaining agents serve as \emph{auditors} who provide evidence about the update's consistency with current memory. Crucially, no agent is assumed honest in our formulation. We define game-theoretical utility functions for both roles and show that the game admits an equilibrium whose payoff structure indicates optimal trust in memory: at equilibrium, a contributor cannot raise an update's influence without also raising its detectability, and an auditor cannot reduce its effort without weakening the consistency signal it produces, which together address the failures as illustrated in Figure \ref{fig:intro} (detailed justification is shown in Appendix \ref{app:failure_modes}).

Guided by this equilibrium, we propose \system, an equilibrium-guided, inference-time calibration mechanism that requires no LLM call for calibration itself. Agents serve only as signal providers: their existing retrieval queries and traversal paths are intercepted and quantified algorithmically against the full memory state. \system instantiates calibration for two representative memory families. For embedding-based memory~\citep{wang2023voyager,zhong2024memorybank}, it checks whether a proposed update disturbs the embedding distribution and reorders retrieval results on auditor queries. For graph-based memory~\citep{anokhin2024arigraph,zhang2025g}, it checks whether new edges are consistent with the local relational neighborhood and with existing multi-hop reasoning paths. In both cases, entries are admitted with trust-discounted weights, so low-confidence content remains available but exerts less influence on downstream retrieval.

We evaluate \system on reasoning- and action-intensive benchmarks~\citep{maharana2024evaluating,shridhar2020alfworld,yang2018hotpotqa,yao2022webshop} across three MAD frameworks~\citep{liu2023dynamic,qian2024scaling,wu2024autogen} and four memory architectures~\citep{anokhin2024arigraph,wang2023voyager,zhang2025g,zhong2024memorybank}. Three findings emerge. (i) \emph{State-aware calibration beats isolated checks:} \system ranks first on every benchmark-framework-memory configuration, confirming that errors missed by per-entry scoring are caught only when updates are evaluated against the global memory state. (ii) \emph{Robustness without trust:} even with up to 50\% adversarial agents, \system stays close to its benign performance, showing that integrity comes from the structural calibration indicator rather than from trusting any individual agent. (iii) \emph{Quality without inflated cost:} since calibration is LLM-free and reuses agents' existing retrieval activity, \system adds zero token overhead and under 5\% latency, showing that strong memory protection is achievable without inflating the inference budget.

In summary, this paper makes the following contributions: \textbf{(1) }We formulate shared-memory multi-agent debate as a zero-trust memory game, in which the equilibrium payoff structure serves as an indicator of optimal memory trust.\textbf{ (2)} Guided by this formulation, we propose \system, a principled mechanism that calibrates memory updates against the global memory state without trusting any individual agent. \textbf{(3) } Through extensive evaluations, \system delivers consistent gains over existing baselines, retains effectiveness under adversarial agents, and incurs negligible inference overhead. Code is released at: \url{https://anonymous.4open.science/r/EquiMem-FE6C}.